\section{Related Work}
\label{sec:relatedwork}

Self-distillation broadly refers to distillation settings where the teacher is
derived from the student model, often with the same architecture or checkpoint.
Our work focuses on recent privileged-context on-policy self-distillation
methods, where the self-teacher is additionally given information unavailable
to the student, such as answers, demonstrations, or feedback. Prior work has
reported gains from \opsd{}
\citep{zhao2026selfdistilledreasoneronpolicyselfdistillation}, SDFT
\citep{shenfeld2026self}, SDPO \citep{hubotter2026reinforcement}, SD-Zero
\citep{he2026selfdistillationzero}, and related recipes, especially for
instruction-tuned models, continual learning, or reasoning compression. Our
result is complementary: we show that in long-budget thinking models,
privileged token-level feedback can degrade the search behavior that enables
test-time reasoning. Concurrent work has linked self-distillation failures to
suppression of epistemic verbalization
\citep{kim2026doesselfdistillationsometimesdegrade}; our analysis studies the
broader fork-suppression mechanism using budget curves, fork/lock diagnostics,
fixed-trajectory token scoring, and trained-student marker shifts. We provide
an expanded discussion of related work in
Appendix~\ref{app:related-work-continued}.
